% =====================================================================
%  PREÁMBULO: paquetes, formato APA 7.ª edición y comandos propios
% =====================================================================
\usepackage[spanish]{babel}
\usepackage[utf8]{inputenc}
\usepackage[T1]{fontenc}
\usepackage{lmodern}
\usepackage[a4paper,top=2.54cm,bottom=2.54cm,left=2.54cm,right=2.54cm,headsep=0.6cm,headheight=15pt]{geometry} % APA 7: márgenes de 2,54 cm

% --- Matemática, figuras y tablas ---
\usepackage{amssymb, amsmath}
\usepackage{graphicx}
\usepackage{float}
\usepackage{subcaption}
\usepackage{booktabs}
\usepackage{threeparttable}
\usepackage{array}
\newcolumntype{L}[1]{>{\raggedright\arraybackslash}p{#1}}
\usepackage{longtable}
\usepackage{pdflscape}
\usepackage{multirow}
\usepackage{pifont}
\usepackage[table]{xcolor}
\usepackage[most]{tcolorbox}
\usepackage{tikz}
\usetikzlibrary{positioning, arrows.meta, babel}
\usepackage{pgfgantt}
\usepackage{algorithm}
\usepackage{algpseudocode}

% --- Texto y listas ---
\usepackage{enumitem}
\usepackage{setspace}
\usepackage{ragged2e}
\usepackage{orcidlink}
\usepackage{url}

% --- Referencias en APA 7.ª edición ---
\usepackage{csquotes}
\usepackage[style=apa,backend=biber,natbib=true]{biblatex}
\addbibresource{bibliografia/referencias.bib}
\setlength{\bibhang}{1.27cm}       % sangría francesa
\setlength{\bibitemsep}{0.5em}

% --- Títulos, índices y número de página ---
\usepackage{titlesec}
\usepackage[titles]{tocloft}
\usepackage{fancyhdr}
\usepackage{hyperref}
\hypersetup{hidelinks, breaklinks=true, pdfstartview=Fit}
\urlstyle{same}

% =====================================================================
%  Colores del título: proceso (azul), tecnología (amarillo), contexto (rojo)
% =====================================================================
\definecolor{procesoAzul}{RGB}{0,82,165}
\definecolor{tecnologiaAmarillo}{RGB}{204,153,0}
\definecolor{contextoRojo}{RGB}{192,0,0}
\newcommand{\proc}[1]{\textcolor{procesoAzul}{#1}}
\newcommand{\tec}[1]{\textcolor{tecnologiaAmarillo}{#1}}
\newcommand{\ctx}[1]{\textcolor{contextoRojo}{#1}}

% colores de los diagramas
\definecolor{diagAzul}{RGB}{58,98,150}
\definecolor{diagVerde}{RGB}{46,125,70}
\definecolor{diagRojo}{RGB}{186,74,52}

% símbolos de la tabla de antecedentes
\newcommand{\cumple}{\textcolor{green!50!black}{\ding{51}}}
\newcommand{\nocumple}{\textcolor{red!70!black}{\ding{55}}}
\newcommand{\parcial}{\textcolor{orange}{$\sim$}}

% =====================================================================
%  Formato APA 7: tablas y figuras
%  Número en negrita, título en cursiva en la línea siguiente, nota debajo
% =====================================================================
\captionsetup[figure]{labelfont=bf, textfont=it, labelsep=newline, name=Figura, justification=raggedright, singlelinecheck=false}
\captionsetup[table]{labelfont=bf, textfont=it, labelsep=newline, name=Tabla, justification=raggedright, singlelinecheck=false}
\counterwithout{figure}{chapter}
\counterwithout{table}{chapter}
\counterwithout{equation}{chapter}

% algoritmos en español
\floatname{algorithm}{Algoritmo}
\algrenewcommand\algorithmicrequire{\textbf{Entrada:}}
\algrenewcommand\algorithmicensure{\textbf{Salida:}}

% =====================================================================
%  APA 7: párrafos con sangría de 1,27 cm y sin espacio adicional
% =====================================================================
\setlength{\parindent}{1.27cm}
\setlength{\parskip}{0pt}
\setlength{\emergencystretch}{3em}

% =====================================================================
%  APA 7: número de página en la esquina superior derecha
% =====================================================================
\pagestyle{fancy}
\fancyhf{}
\fancyhead[R]{\thepage}
\renewcommand{\headrulewidth}{0pt}
\fancypagestyle{plain}{\fancyhf{}\fancyhead[R]{\thepage}\renewcommand{\headrulewidth}{0pt}}

% =====================================================================
%  Títulos: "1. Aspectos Generales", "1.1. ...", "1.1.1. ..."
% =====================================================================
\titleformat{\chapter}[block]{\normalfont\bfseries\Large\centering}{\thechapter.}{0.5em}{}
\titlespacing*{\chapter}{0pt}{0pt}{1.2cm}
\titleformat{name=\chapter,numberless}[block]{\normalfont\bfseries\Large\centering}{}{0pt}{}
\titlespacing*{name=\chapter,numberless}{0pt}{0pt}{1cm}
\titleformat{\section}{\normalfont\bfseries\large}{\thesection.}{0.6em}{}
\titleformat{\subsection}{\normalfont\bfseries\normalsize}{\thesubsection.}{0.6em}{}
\titleformat{\subsubsection}{\normalfont\bfseries\itshape\normalsize}{}{0pt}{}
\setcounter{secnumdepth}{2}
\setcounter{tocdepth}{2}

% índice general como el modelo de la escuela
\addto\captionsspanish{%
  \renewcommand{\contentsname}{Índice General}%
  \renewcommand{\listtablename}{Índice de Tablas}%
  \renewcommand{\listfigurename}{Índice de Figuras}%
  \renewcommand{\tablename}{Tabla}%
}
\renewcommand{\cftchapaftersnum}{.}
\renewcommand{\cftsecaftersnum}{.}
\renewcommand{\cftsubsecaftersnum}{.}
\renewcommand{\cftchapleader}{\hfill}
\setlength{\cftsecnumwidth}{2.6em}
\setlength{\cftsubsecnumwidth}{3.4em}
\renewcommand{\cfttoctitlefont}{\hfill\normalfont\Large\bfseries}
\renewcommand{\cftaftertoctitle}{\hfill}
\renewcommand{\cftloftitlefont}{\hfill\normalfont\Large\bfseries}
\renewcommand{\cftafterloftitle}{\hfill}
\renewcommand{\cftlottitlefont}{\hfill\normalfont\Large\bfseries}
\renewcommand{\cftafterlottitle}{\hfill}

% índice de algoritmos
\renewcommand{\listalgorithmname}{Índice de Algoritmos}

% índice de ecuaciones
\newcommand{\listequationsname}{Índice de Ecuaciones}
\newlistof{ecuaciones}{equ}{\listequationsname}
\newcommand{\ecuacion}[1]{\addcontentsline{equ}{ecuaciones}{\protect\numberline{\theequation}#1}}
\setlength{\cftecuacionesnumwidth}{2.5em}
\renewcommand{\cftequtitlefont}{\hfill\normalfont\Large\bfseries}
\renewcommand{\cftafterequtitle}{\hfill}

% índice de anexos
\newcommand{\listanexosname}{Índice de Anexos}
\newlistof{anexos}{anx}{\listanexosname}
\setlength{\cftanexosnumwidth}{2.5em}
\renewcommand{\cftanxtitlefont}{\hfill\normalfont\Large\bfseries}
\renewcommand{\cftafteranxtitle}{\hfill}

% anexos: aparecen en el índice general y en el índice de anexos
\newcommand{\registraranexo}[2]{%
  \addcontentsline{toc}{chapter}{Anexo #1. #2}%
  \addcontentsline{anx}{anexos}{\protect\numberline{#1}#2}}
\newcommand{\anexo}[2]{%
  \clearpage
  \phantomsection
  \registraranexo{#1}{#2}%
  {\noindent\large\bfseries Anexo #1. #2\par}\vspace{0.6cm}}
